% SGA 6, Exposición VI: fragmento español alineado con la fuente.
% Alcance: sección 1, números 1.8--1.10, líneas 8691--8783 inclusive.
% Fuente francesa: sga6_fr_r1/src/fr.tex.
% SHA-256 de la fuente francesa completa:
% 0D2158B9FE3B466309A92165727A0B191422E05FBAA0DACD5806E8322C54BA1A.
\providecommand{\sisi}[2]{#2}

\textbf{1.8.} Volviendo a la demostración del teorema 1.1, mostremos que \(K^\circ(X)\) está engendrado como \(K^\circ(S)\)-módulo por las potencias de \(\xi\).

Si \(F\) es un \(\mathcal O_X\)-módulo localmente libre de tipo finito, existe, según 1.4, un \(B\)-módulo graduado de presentación finita \(M\), plano sobre \(S\), tal que \(F=\Projz(M)\). Según 1.6, \(M\) admite una resolución graduada por \(B\)-módulos graduados localmente libres de tipo finito. Se deduce una sucesión exacta:
\[
  0\longrightarrow \Projz(L_{r+1})\longrightarrow\cdots\longrightarrow
  \Projz(L_0)\longrightarrow\Projz(M)\longrightarrow 0,
\]
y la igualdad siguiente en \(K^\circ(X)\):
\[
  \cl(F)=\sum_i(-1)^i\cl\bigl(\Projz(L_i)\bigr).
\]
Utilizando la filtración sobre los \(L_i\) definida en 1.7, se obtiene:
\[
  \cl(\Projz(L_i))=\sum_j \cl\bigl(\Projz(N_{ij}\otimes_{\mathcal O_S}B(j))\bigr).
\]
Ahora bien, \(\Projz(N_{ij}\otimes_{\mathcal O_S}B(j))=f^*(N_{ij})(j)\); se obtiene, pues:
\[
\begin{aligned}
  \cl(\Projz(L_i))
  &=\sum_j \cl\bigl(f^*(N_{ij})\bigr).\xi^j \\
  &=\sum_j f^*\bigl(\cl(N_{ij})\bigr).\xi^j.
\end{aligned}
\]
Se ha expresado finalmente \(\cl(F)\) como combinación lineal, con coeficientes en \(K^\circ(S)\), de las potencias de \(\xi\).

\textbf{1.9.} Para ver que \(1,\xi,\ldots,\xi^{r+1}\) son linealmente dependientes sobre \(K^\circ(S)\), se considera el homomorfismo canónico sobreyectivo \(f^*(E)\to\mathcal O_X(1)\), cuyo núcleo \(F\) es localmente libre de rango \(r\). Se puede, pues, escribir:
\[
  \lambda^{r+1}\bigl(f^*(E)-\xi\bigr)=0,
  \qquad \text{con } E=\cl E,
\]
estando definidas las operaciones \(\lambda^i\) sobre \(K^\circ(X)\) en V 2.2. Se obtiene, pues:
\[
  \sum_{i=0}^{r+1}\lambda^i(f^*(E))\lambda^{r+1-i}(-\xi)=0.
\]
Como \(\mathcal O_X(1)\) es un haz invertible, \(\lambda_t(\xi)=1+\xi t\), y por consiguiente \(\lambda_t(-\xi)=1/(1+\xi t)\), de donde resulta finalmente la relación:
\begin{equation}\tag{1.9.1}
  \sum_{i=0}^{r+1}(-1)^{r+1-i}\xi^{r+1-i}\lambda^i(E)=0.
\end{equation}

Mostremos finalmente que \(1,\xi,\ldots,\xi^r\) son independientes sobre \(K^\circ(S)\)\footnote{Esta demostración se debe a J. P. Jouanolou.}.


\textbf{Lema 1.10.} Con las notaciones y las hipótesis de 1.1, sean \(F\) y \(G\) dos haces localmente libres de tipo finito sobre \(X\), tales que:
\[
  \cl_{K^\circ(X)}(F)=\cl_{K^\circ(X)}(G);
\]
entonces existe un entero \(n_0\) tal que, para \(n\geq n_0\), se tiene:
\[
\begin{array}{ll}
\text{i)}& f_*(F(n)),\ f_*(G(n))\text{ son localmente libres de tipo finito sobre }S;\\[1mm]
\text{ii)}& \cl_{K^\circ(S)}\bigl(f_*(F(n))\bigr)=\cl_{K^\circ(S)}\bigl(f_*(G(n))\bigr).
\end{array}
\]

Sea \(A\) el grupo abeliano libre engendrado por las clases (salvo isomorfismo) de haces localmente libres de tipo finito sobre \(X\). Decir que, en \(K^\circ(X)\), \(F\) y \(G\) tienen la misma clase significa, según la construcción de \(K^\circ\), que existen sucesiones exactas de \(\mathcal O_X\)-módulos localmente libres de tipo finito:
\[
  0\longrightarrow L_i\longrightarrow M_i\longrightarrow N_i\longrightarrow 0,
  \qquad i=1,\ldots,k,
\]
y enteros \(n_i\), tales que en \(A\) se tiene la igualdad:
\[
  (F)-(G)=\sum_i n_i\bigl((L_i)-(M_i)+(N_i)\bigr).
\]
Para todo \(n\), la tensorización por \(\mathcal O_X(n)\) define un automorfismo de \(A\), de donde:
\[
  (F(n))-(G(n))=\sum_i \sisi{}{n_i}\bigl((L_i(n))-(M_i(n))+(N_i(n))\bigr),
\]
y las sucesiones \(0\to L_i(n)\to M_i(n)\to N_i(n)\to0\) siguen siendo exactas.

Sean \(B\) el grupo abeliano libre engendrado por las clases (salvo isomorfismo) de haces cuasicoherentes sobre \(S\), y \(A'\subset B\) el subgrupo engendrado por las clases de haces localmente libres de tipo finito sobre \(S\). Mediante imagen directa se define un homomorfismo \(A\to B\), y se tiene aún:
\[
\begin{aligned}
  \bigl(f_*(F(n))\bigr)-\bigl(f_*(G(n))\bigr)
  ={}&\sum_i n_i\bigl((f_*(L_i(n)))-(f_*(M_i(n)))\\
      &\hspace{35mm}+(f_*(N_i(n)))\bigr).
\end{aligned}
\]
Según \sisi{1.2}{1.3}, para \(n\) suficientemente grande se cumplirán las condiciones siguientes:
\[
  \forall k\geq1,\ \forall i,
  \quad R^k f_*(L_i(n))=0;
\]
\[
  \forall i,
  \quad f_*(L_i(n)),\ f_*(M_i(n)),\ f_*(N_i(n)),\ f_*(F(n)),\ f_*(G(n))
  \text{ son localmente libres de tipo finito.}
\]
De ello resulta que, para \(n\) suficientemente grande, todos estos elementos pertenecen a \(A'\), verifican la relación escrita anteriormente, y las sucesiones
\[
  0\longrightarrow f_*(L_i(n))\longrightarrow f_*(M_i(n))\longrightarrow f_*(N_i(n))\longrightarrow0
\]
son exactas. Esto implica que, en \(K^\circ(S)\), \(f_*(F(n))\) y \(f_*(G(n))\) tienen la misma clase.

